% GENERATED FILE. Edit canonical lessons, then run npm run generate:books.
% canonical-source-hash: fnv1a64:f52ea4d567a59077
% canonical-lessons: BN-C456-choritro, BN-C456-gun, BN-C456-obhyas, BN-C456-achoron, BN-C456-ropto-kora

\chapter{A Person's Character}
\label{ch:bn-a-person-s-character}
\hlchaptermodality{456}

\begin{chapteropening}
\textbf{By the end of this chapter:} \emph{I can describe someone's character and habits.}

Complete the last lesson of chapter 456: I can describe someone's character and habits.
\end{chapteropening}

\section[\texorpdfstring{\bn{চরিত্র} (chôritrô)}{চরিত্র (chôritrô)}]{\bn{চরিত্র} (chôritrô) — character; moral nature}

\label{lesson:BN-C456-choritro}



\begin{quote}
Before the new one: say the Bengali for eager, then the Bengali for to be able to wait.
\end{quote}

\subsection*{You'll want to know: \bn{চরিত্র}}
\textbf{\bn{চরিত্র}} — \emph{chôritrô} — \textquotedblleft{}character; moral nature\textquotedblright{}.

\begin{quote}
\bn{লোকটার} \bn{চরিত্র} \bn{ভালো} \bn{নয়।}
\end{quote}

\begin{quote}
\emph{lokṭār chôritrô bhālo nôy.}
\end{quote}

\begin{quote}
That man does not have a good character.
\end{quote}

\bn{চরিত্র} is also a character in a story: \bn{গল্পের} \bn{চরিত্র}.

\begin{cousinweb}
\textbf{\bn{চরিত্র}} (\emph{chôritrô}) is a tatsama, Sanskrit \emph{caritra}, \textquotedblleft{}going\textquotedblright{}, then \textquotedblleft{}behaviour, deeds\textquotedblright{} and \textquotedblleft{}nature, disposition\textquotedblright{} (Monier-Williams).
\end{cousinweb}

\begin{grammarlens}[title={What you've built}]
One more word for describing what you lived and felt.
\end{grammarlens}

\subsection*{Your turn}
\begin{itemize}
\raggedright
  \item \emph{Say it:} \emph{chôritrô}
  \item \emph{Say it:} \emph{chôritrô}, once more
  \item \emph{Say it:} \emph{lokṭār chôritrô bhālo nôy.} — That man does not have a good character.
  \item \emph{Recall:} say the Bengali for eager, then the Bengali for to be able to wait, then say \emph{chôritrô} again
\end{itemize}

\begin{tcolorbox}[breakable,colback=teal!4,colframe=teal!35!black,title={Before you move on}]
What does \bn{চরিত্র} mean? (\textquotedblleft{}Character; moral nature\textquotedblright{}.) Say it once more.
\end{tcolorbox}
\section[\texorpdfstring{\bn{গুণ} (gun)}{গুণ (gun)}]{\bn{গুণ} (gun) — a virtue; a good quality}

\label{lesson:BN-C456-gun}



\begin{quote}
Before the new one: say the Bengali for to be able to wait, then the Bengali for character.
\end{quote}

\subsection*{You'll want to know: \bn{গুণ}}
\textbf{\bn{গুণ}} — \emph{gun} — \textquotedblleft{}a virtue; a good quality\textquotedblright{}.

\begin{quote}
\bn{ধৈর্য} \bn{মীরার} \bn{সবচেয়ে} \bn{বড়} \bn{গুণ।}
\end{quote}

\begin{quote}
\emph{dhoirjô mīrār sôbcheye bôṛo gun.}
\end{quote}

\begin{quote}
Patience is Mira's greatest virtue.
\end{quote}

\begin{cousinweb}
\textbf{\bn{গুণ}} (\emph{gun}) is a tatsama, Sanskrit \emph{guṇa}, first \textquotedblleft{}a single strand of a cord\textquotedblright{} and then \textquotedblleft{}a quality, a virtue\textquotedblright{} (Monier-Williams).
\end{cousinweb}

\begin{grammarlens}[title={What you've built}]
One more word for describing what you lived and felt.
\end{grammarlens}

\subsection*{Your turn}
\begin{itemize}
\raggedright
  \item \emph{Say it:} \emph{gun}
  \item \emph{Say it:} \emph{gun}, once more
  \item \emph{Say it:} \emph{dhoirjô mīrār sôbcheye bôṛo gun.} — Patience is Mira's greatest virtue.
  \item \emph{Recall:} say the Bengali for to be able to wait, then the Bengali for character, then say \emph{gun} again
\end{itemize}

\begin{tcolorbox}[breakable,colback=teal!4,colframe=teal!35!black,title={Before you move on}]
What does \bn{গুণ} mean? (\textquotedblleft{}A virtue; a good quality\textquotedblright{}.) Say it once more.
\end{tcolorbox}
\section[\texorpdfstring{\bn{অভ্যাস} (ôbhyās)}{অভ্যাস (ôbhyās)}]{\bn{অভ্যাস} (ôbhyās) — a habit}

\label{lesson:BN-C456-obhyas}



\begin{quote}
Before the new one: say the Bengali for character, then the Bengali for a good quality.
\end{quote}

\subsection*{You'll want to know: \bn{অভ্যাস}}
\textbf{\bn{অভ্যাস}} — \emph{ôbhyās} — \textquotedblleft{}a habit\textquotedblright{}.

\begin{quote}
\bn{রোজ} \bn{ভোরে} \bn{হাঁটা} \bn{দাদুর} \bn{পুরনো} \bn{অভ্যাস।}
\end{quote}

\begin{quote}
\emph{roj bhore hā̃ṭā dādur purôno ôbhyās.}
\end{quote}

\begin{quote}
Walking at dawn every day is an old habit of my grandfather's.
\end{quote}

\begin{cousinweb}
\textbf{\bn{অভ্যাস}} (\emph{ôbhyās}) is a tatsama, Sanskrit \emph{abhyāsa}, \textquotedblleft{}repetition, practice, habit\textquotedblright{} (Monier-Williams): what you do again and again.
\end{cousinweb}

\begin{grammarlens}[title={What you've built}]
One more word for describing what you lived and felt.
\end{grammarlens}

\subsection*{Your turn}
\begin{itemize}
\raggedright
  \item \emph{Say it:} \emph{ôbhyās}
  \item \emph{Say it:} \emph{ôbhyās}, once more
  \item \emph{Say it:} \emph{roj bhore hā̃ṭā dādur purôno ôbhyās.} — Walking at dawn every day is an old habit of my grandfather's.
  \item \emph{Recall:} say the Bengali for character, then the Bengali for a good quality, then say \emph{ôbhyās} again
\end{itemize}

\begin{tcolorbox}[breakable,colback=teal!4,colframe=teal!35!black,title={Before you move on}]
What does \bn{অভ্যাস} mean? (\textquotedblleft{}A habit\textquotedblright{}.) Say it once more.
\end{tcolorbox}
\section[\texorpdfstring{\bn{আচরণ} (āchôrôn)}{আচরণ (āchôrôn)}]{\bn{আচরণ} (āchôrôn) — conduct; the way someone behaves}

\label{lesson:BN-C456-achoron}



\begin{quote}
Before the new one: say the Bengali for a good quality, then the Bengali for a habit.
\end{quote}

\subsection*{You'll want to know: \bn{আচরণ}}
\textbf{\bn{আচরণ}} — \emph{āchôrôn} — \textquotedblleft{}conduct; the way someone behaves\textquotedblright{}.

\begin{quote}
\bn{ছেলেটির} \bn{ভদ্র} \bn{আচরণ} \bn{আমার} \bn{ভালো} \bn{লাগল।}
\end{quote}

\begin{quote}
\emph{chheleṭir bhôdrô āchôrôn āmār bhālo lāglô.}
\end{quote}

\begin{quote}
I liked the boy's polite conduct.
\end{quote}

\begin{cousinweb}
\textbf{\bn{আচরণ}} (\emph{āchôrôn}) is a tatsama, Sanskrit \emph{ācaraṇa}, \textquotedblleft{}undertaking, practising, conduct\textquotedblright{} (Monier-Williams).
\end{cousinweb}

\begin{grammarlens}[title={What you've built}]
One more word for describing what you lived and felt.
\end{grammarlens}

\subsection*{Your turn}
\begin{itemize}
\raggedright
  \item \emph{Say it:} \emph{āchôrôn}
  \item \emph{Say it:} \emph{āchôrôn}, once more
  \item \emph{Say it:} \emph{chheleṭir bhôdrô āchôrôn āmār bhālo lāglô.} — I liked the boy's polite conduct.
  \item \emph{Recall:} say the Bengali for a good quality, then the Bengali for a habit, then say \emph{āchôrôn} again
\end{itemize}

\begin{tcolorbox}[breakable,colback=teal!4,colframe=teal!35!black,title={Before you move on}]
What does \bn{আচরণ} mean? (\textquotedblleft{}Conduct; the way someone behaves\textquotedblright{}.) Say it once more.
\end{tcolorbox}
\section[\texorpdfstring{\bn{রপ্ত} \bn{করা} (rôptô kôrā)}{রপ্ত করা (rôptô kôrā)}]{\bn{রপ্ত} \bn{করা} (rôptô kôrā) — to master (a skill); to get the hang of}

\label{lesson:BN-C456-ropto-kora}



\begin{quote}
Before the new one: say the Bengali for a habit, then the Bengali for conduct.
\end{quote}

\subsection*{You'll want to know: \bn{রপ্ত} \bn{করা}}
\textbf{\bn{রপ্ত} \bn{করা}} — \emph{rôptô kôrā} — \textquotedblleft{}to master (a skill); to get the hang of\textquotedblright{}.

\begin{quote}
\bn{অরুণ} \bn{তিন} \bn{মাসে} \bn{সাঁতার} \bn{রপ্ত} \bn{করল।}
\end{quote}

\begin{quote}
\emph{ôrun tin māse sā̃tār rôptô kôrlô.}
\end{quote}

\begin{quote}
Arun mastered swimming in three months.
\end{quote}

Past (he/she …): \bn{তিনি} \bn{কাজটা} \bn{রপ্ত} \bn{করলেন}, \emph{tini kājṭā rôptô kôrlen}.

\begin{cousinweb}
\textbf{\bn{রপ্ত} \bn{করা}} (\emph{rôptô kôrā}) is \emph{rôptô}, \textquotedblleft{}practised, mastered\textquotedblright{}, whose origin is not settled, with \emph{kôrā}, \textquotedblleft{}to do\textquotedblright{} (CDIAL 2814).
\end{cousinweb}

\begin{grammarlens}[title={What you've built}]
One more word for describing what you lived and felt.
\end{grammarlens}

\subsection*{Your turn}
\begin{itemize}
\raggedright
  \item \emph{Say it:} \emph{rôptô kôrā}
  \item \emph{Say it:} \emph{rôptô kôrā}, once more
  \item \emph{Say it:} \emph{ôrun tin māse sā̃tār rôptô kôrlô.} — Arun mastered swimming in three months.
  \item \emph{Recall:} say the Bengali for a habit, then the Bengali for conduct, then say \emph{rôptô kôrā} again
\end{itemize}

\begin{tcolorbox}[breakable,colback=teal!4,colframe=teal!35!black,title={Before you move on}]
What does \bn{রপ্ত} \bn{করা} mean? (\textquotedblleft{}To master (a skill); to get the hang of\textquotedblright{}.) Say it once more.
\end{tcolorbox}
